\icode{N14H-MV-SI-L2-261011-001}
\nceeID[6]{2014安徽}{14}
在科学研究中，科学家常将未知现象同已知现象进行比较，找出其共同点，进一步推测未知现象的特性和规律。法国物理学家库伦在研究异种电荷的吸引问题时，曾将扭秤的振动周期与电荷间距离的关系类比单摆的振动周期与摆球到地心距离的关系。已知单摆摆长为\(l\)，引力常量为\(G\)，地球质量为\(M\)，摆球到地心的距离为\(r\)，则单摆振动周期\(T\)与距离\(r\)的关系式为\paren[]
\begin{choices}
	\item \(T = 2\pi r \sqrt{\frac{GM}{l}}\)\
	\item \(T = 2\pi r \sqrt{\frac{l}{GM}}\)
	\item \(T = \frac{2\pi}{r}\sqrt{\frac{GM}{l}}\)
	\item \(T = 2\pi l \sqrt{\frac{l}{GM}}\)
\end{choices}
\begin{answer}
	B\qquad 提示：\(T = 2\pi \sqrt{\frac{l}{g}}\)\quad \(g = \frac{GM}{r^2}\)
\end{answer}